\subsection{The additional unit condition in the quintic capitulation criterion}

\begin{proposition}[KT-ORDER-FIVE-GREENBERG-UNIT-GATE]
Let $p\equiv1\pmod {10}$ be prime, let $F$ be the cyclic quintic
field of conductor $p$, and let $F_n/F$ be the degree-$p^n$ layer
of its cyclotomic $\mathbf Z_p$-extension.  Put
\[
 A_n=\operatorname{Cl}(F_n)\{p\},\qquad
 E_n=\mathcal O_{F_n}^{\times},\qquad
 G_n=\operatorname{Gal}(F_n/F),
 \qquad j_n:A_0\longrightarrow A_n.
                                                        \tag{O5G.1}
\]
For every $n\ge1$, all units in $E_0$ are norms from
$F_n^{\times}$, and there is an exact sequence
\[
 1\longrightarrow j_n(A_0)\longrightarrow A_n^{G_n}
 \longrightarrow E_0/N_{F_n/F}E_n\longrightarrow1,
 \qquad |A_n^{G_n}|=|A_0|.                              \tag{O5G.2}
\]
Consequently
\[
 [E_0:N_{F_n/F}E_n]=|\ker(j_n)|.                        \tag{O5G.3}
\]
If the remaining possibility in KT-ORDER-FIVE-LENGTH-ONE is
$A_0\simeq\mathbf Z/p\mathbf Z$, then the following are equivalent:
\begin{enumerate}
\item $[E_0:N_{F_n/F}E_n]=p$ for some $n$;
\item $A_0$ capitulates in some $F_n$;
\item $\varinjlim_n A_n=0$, where the transition maps extend ideals;
\item the orders $|A_n|$ are bounded, equivalently the cyclotomic
      Iwasawa invariants satisfy $\lambda=\mu=0$.
\end{enumerate}
In particular, the required extra condition concerns norms of
\emph{units}; the assertion that every base unit is a norm of a
field element already holds without capitulation.
\end{proposition}

\begin{proof}
The field $F_n$ is the unique degree-$5p^n$ subfield of
$\mathbf Q(\mu_{p^{n+1}})$.  It is totally real and has a unique
prime $\mathfrak p_n$ above $p$, totally ramified over $\mathbf Q$.
Moreover,
\[
 \pi_n=N_{\mathbf Q(\mu_{p^{n+1}})/F_n}
                       (1-\zeta_{p^{n+1}}),\qquad
                    (\pi_n)=\mathfrak p_n.              \tag{O5G.4}
\]
Indeed the ideal norm of the unique prime in the full cyclotomic
field has exponent its residue degree, which is one.  Thus the
only prime ramified in $F_n/F$ has principal ideal class.

At every finite place other than $p$, a base unit is a local norm:
the local extension is unramified and the norm on local units is
surjective.  At real places there is no complexification.
The product formula for the cyclic norm-residue symbols then
forces the symbol at the unique $p$-place to be trivial as well.
The Hasse norm theorem for the cyclic extension $F_n/F$ gives
\[
                    E_0\subset N_{F_n/F}(F_n^{\times}). \tag{O5G.5}
\]

Here is a direct class-field proof of the order assertion in
(O5G.2).  Let $H_n$ be the ordinary Hilbert $p$-class field of
$F_n$, and let $M_n$ be its maximal subfield abelian over $F$.
The field $H_n$ is Galois over $F$, and
\[
 \operatorname{Gal}(M_n/F_n)
       \simeq A_n/(\sigma-1)A_n,                       \tag{O5G.6}
\]
for a generator $\sigma$ of $G_n$.
This follows because the commutator subgroup is
$(\sigma-1)A_n$: after quotienting by this subgroup the kernel
is central and the quotient is cyclic, so the entire group is
abelian.  In the abelian extension $M_n/F$, the inertia group
$I$ at $p$ has order exactly $p^n$, since $M_n/F_n$ is
unramified and $F_n/F$ is totally ramified.  It maps
isomorphically to $G_n$.  Its fixed field $M_n^I$ is everywhere
unramified over $F$, including the real places, and hence is
contained in $H_0$.  Thus $M_n\subset F_nH_0$.
Conversely $F_nH_0/F_n$ is unramified abelian of $p$-power degree
and $F_nH_0/F$ is abelian, so $F_nH_0\subset M_n$.
Total ramification gives $F_n\cap H_0=F$, and therefore
$|A_n/(\sigma-1)A_n|=|A_0|$.
For a finite module the kernel and cokernel of $\sigma-1$
have equal orders, proving $|A_n^{G_n}|=|A_0|$.

For completeness, the ambiguous-class map sends an invariant
class represented by $\mathfrak a$, with
$\mathfrak a^{\sigma-1}=(x)$, to $N_{F_n/F}x$ modulo
$N_{F_n/F}E_n$.  This is well-defined.  Its kernel consists
of classes with an invariant ideal representative: if the
norm equals the norm of a unit, divide $x$ by that unit and
apply Hilbert's theorem 90.  Conversely an invariant ideal
representative gives the trivial unit class.
The map is onto $E_0\cap N(F_n^{\times})/N E_n$:
for a unit $N x$, the ideal $(x)$ has ideal norm one, and the
norm-one kernel on the permutation module of fractional ideals
is $(\sigma-1)I(F_n)$.
An invariant fractional ideal is a product of an extended ideal
of $F$ and powers of $\mathfrak p_n$.  Equation (O5G.4) removes
the latter classes.  Taking $p$-primary parts and using (O5G.5)
therefore gives exactly (O5G.2).  The unit quotient is a
$p$-group, since the $p^n$-th power of any base unit is its norm
when regarded as a unit upstairs.  Orders now prove (O5G.3).

Under $|A_0|=p$, (O5G.3) proves the equivalence of the first
two assertions.  The unique $p$-prime is totally ramified in
the whole tower, so Greenberg's one-prime capitulation criterion
applies.  In the notation of Ichimura's published version,
Theorem~1 gives $\varinjlim A_n=0$ if the image of $A_0$ is zero;
the converse is immediate from finiteness of $A_0$.
Under the same one-prime hypothesis, his Remark~1(1) identifies
the auxiliary extension $L_n$ with $H_nF_\infty$.
Total ramification gives $H_n\cap F_\infty=F_n$, so
$[L_n:F_\infty]=|A_n|$.  Remark~1(2) identifies vanishing of
the direct limit with finiteness of $\bigcup_nL_n/F_\infty$,
which is therefore equivalent to bounded $|A_n|$.
The standard Iwasawa class-number formula gives the stated
equivalence with $\lambda=\mu=0$.
\end{proof}

\paragraph{What the criterion does and does not supply.}
When $A_0\simeq\mathbf Z/p\mathbf Z$, the norm-unit index is
either $1$ or $p$.  If it is $1$ at every layer, $j_n$ is
injective at every layer.  If it becomes $p$, the base class
capitulates and the tower has bounded class numbers.
The order-five class bound does not evaluate this additional
norm-unit index.  The cited criterion assumes capitulation,
equivalently the displayed norm-unit condition; it does not
deduce it from cyclicity of $A_0$.
Even the bounded-growth conclusion asserts capitulation of
classes after extension, not their triviality in $F$.
No order-five component is removed by this proposition.

\paragraph{Exact primary scope.}
H.~Ichimura, \emph{A note on the ideal class group of the
cyclotomic $\mathbf Z_p$-extension of a totally real number
field}, Acta Arithmetica \textbf{105.1} (2002), 29--34:
Theorems~1--2 and Proposition~1 are on p.~29;
Remark~1(2) and the imported bounded-unit-index Lemma~1 are on
p.~30; the complete proofs of Lemma~2, Theorem~2 and
Proposition~1 are on p.~31.  These pages were read in full.
This is a direct published exposition of Greenberg's
1976 Theorem~1; no fresh read of the original 1976 proof is
claimed.  The norm-unit exact sequence is also stated as
equation (1.1), p.~4, of G.~Gras,
\emph{Algebraic norms and capitulation of $p$-class groups in
ramified cyclic $p$-extensions}, arXiv:2211.12279v3
(6 February 2023).  Its general setup on pp.~2--5 and
cyclotomic discussion on pp.~23--25 were read.
The restricted tame setting of that paper's Section~2.2 is
not substituted for the present cyclotomic extension.
The argument above proves the specialized exact sequence and
ambiguous-class order directly; the Hasse norm theorem,
Hilbert's theorem 90 and global class field theory are imported
foundational results, not freshly reproved here.
